\section{E0: does layer 1 produce a usable representation?}
\label{sec:density}

Before comparing architectures it is worth asking whether the intermediate representation
can support a second layer at all. A layer-2 clause is a conjunction over a patch of layer-1
bits; if those bits fire with probability $p$, a conjunction of $j$ of them is satisfied with
probability $\approx p^{\,j}$. The sweep trains layer~1 at several specificities and
clause-size budgets on \densitySubset{} images for \densityEpochs{} epochs and measures the
firing rate of the resulting maps.

\begin{table}[h]
\centering\small
\caption{Layer-1 clause density against specificity. \code{fire} is the median per-channel
firing rate at OR-pool $1$, $2$ and $4$; \code{dead} is the fraction of channels that never
fire. Rows with a clause-size budget $b$ are diagnostics of the budget mechanism
(\S\ref{sec:budgetbug}); the sequential row trains on a tenth of the data, so its accuracy is
not comparable.}
\label{tab:density}
\pgfplotstabletypeset[
  col sep=comma, string type,
  every head row/.style={before row=\ttoprule, after row=\ttmidrule},
  every last row/.style={after row=\ttbotrule},
  columns={label,includes,fire1,fire2,fire4,dead2,l1_acc},
  columns/label/.style={column name=\thead{layer-1 setting}, column type={l}},
  columns/includes/.style={column name=\thead{includes}, column type={r}},
  columns/fire1/.style={column name=\thead{fire (p1)}, column type={r}},
  columns/fire2/.style={column name=\thead{fire (p2)}, column type={r}},
  columns/fire4/.style={column name=\thead{fire (p4)}, column type={r}},
  columns/dead2/.style={column name=\thead{dead}, column type={r}},
  columns/l1_acc/.style={column name=\thead{L1 acc (\%)}, column type={r}},
]{data/density.csv}
\end{table}

\begin{figure}[h]
\centering
\begin{tikzpicture}
\begin{axis}[
  width=0.92\textwidth, height=6.2cm,
  xlabel={clause-activation channel (sorted)}, ylabel={firing rate},
  ymode=log, ymin=1e-6, ymax=1,
  grid=major, grid style={ttline!40},
  label style={font=\sffamily\small}, tick label style={font=\sffamily\footnotesize},
  every axis plot/.append style={very thick},
]
\addplot[ttorange] table[col sep=comma, x=channel, y=rate] {data/firing_curve.csv};
\addplot[ttslate, dashed, domain=0:\cOne] {0.05};
\legend{\densityDensestLabel{} (densest setting), {usable band ($5\%$)}}
\end{axis}
\end{tikzpicture}
\caption{Per-channel firing rate of the layer-1 clause maps after $2\times2$ OR-pooling, for
the densest layer-1 setting found (\densityDensestLabel{}), sorted. Log scale. The dashed
line marks the density a layer-2 conjunction would need to have a realistic chance of firing.}
\label{fig:firing}
\end{figure}

\subsection{Result}

The densest layer-1 setting the sweep can reach (\densityDensestLabel{}) gives a median
per-channel firing rate of \densityDensestFireRaw{}\% raw, \densityDensestFire{}\% after
$2\times2$ OR-pooling and \densityDensestFirePoolFour{}\% after $4\times4$, with a median
clause size of \densityDensestIncludes{} literals. At the other end,
\densitySparsestLabel{} yields \densitySparsestFire{}\% with
\densitySparsestDead{}\% of channels dead.

\begin{ttkey}
Even in its densest configuration the intermediate representation is far sparser than a
conjunction over it can exploit. A two-literal layer-2 clause drawing from channels that fire
at \densityDensestFire{}\% matches roughly
$(\densityDensestFire{}/100)^2$ of positions. OR-pooling is the only lever that moves this
materially, which is why it is treated as load-bearing rather than optional.
\end{ttkey}

\subsection{An interaction worth reporting: the clause-size budget stops binding}
\label{sec:budgetbug}

\code{max\_included\_literals} is documented as a clause-size budget: a matching clause over
budget receives no Type~Ia feedback, and its Type~I events act as Type~Ib until it shrinks
back. It gates Type~Ia only. Type~II is ungated --- it pushes \emph{excluded} literals
towards inclusion, one step per event, capped only by the distance to the include boundary:

\begin{lstlisting}[language=Python]
inc2 = torch.minimum(n2, room.to(n2.dtype), out=n2)   # functional.py
\end{lstlisting}

Under \code{feedback\_mode="batch"} a single commit aggregates the Type~II events of the
whole mini-batch, and a convolutional model contributes one opportunity per patch --- $841$
for $4\times4$ patches on a $32\times32$ image. Enough events accumulate in one commit to
carry many literals across the boundary at once, so the budget never binds. Measured at
budget $b=8$: median clause size \incBudgetBatchFifty{} at batch~50,
\incBudgetBatchFive{} at batch~5, and \incBudgetSeq{} under sequential feedback.

This is a fidelity cost of batched feedback beyond the one the library already documents, and
it is specific to the convolutional models, where the patch dimension multiplies the event
count. For the experiments below, specificity $s$ is therefore used as the density knob
rather than the clause-size budget.
